% Beamer narrative — STR prohibitions and tract rents (Callaway & Sant'Anna spine)
\documentclass[11pt,aspectratio=169]{beamer}

\usepackage[T1]{fontenc}
\usepackage[utf8]{inputenc}
\usepackage{lmodern}
\usepackage{graphicx}
\usepackage{booktabs}
\usepackage{array}
\usepackage{xcolor}
\usepackage{tikz}

\usetheme{UChicago}

% Pipeline outputs resolved from docs/white-paper/ppt -> repo root ../../../
\newcommand{\WpOut}{../../../output/did-cs-whitepaper}
\graphicspath{{../../../output/did-cs-whitepaper/}{../../../output/did-cs/}}

% Image with graceful placeholder if the figure has not been built yet.
\newcommand{\WPFig}[2]{%
  \IfFileExists{\WpOut/#1}{%
    \includegraphics[#2]{#1}%
  }{%
    \fbox{\parbox{0.9\linewidth}{\centering\scriptsize Missing file: \detokenize{#1}\\[0.35em]Run \texttt{make run-did-pipeline-cs-whitepaper}}}%
  }%
}

% Snippet \input with placeholder.
\newcommand{\WPInput}[1]{%
  \IfFileExists{\WpOut/#1}{\input{\WpOut/#1}}%
  {{\scriptsize\textit{Missing snippet: \texttt{\detokenize{#1}}}.}}%
}

\title[STR prohibitions and rents]{Short-Term Rental Prohibitions and Census-Tract Rents}
\subtitle{A staggered difference-in-differences design with cohort-aware estimation}
\author[DSI Housing]{Data Science Institute Housing Clinic \and City of Chicago DTI}
\date{\today}

\begin{document}

% --------------------------------------------------------------------
\begin{frame}[plain]
  \titlepage
\end{frame}

% --------------------------------------------------------------------
\begin{frame}{What we ask, and why it is hard}
\begin{itemize}
  \item \textbf{Question.} Do census-tract STR prohibitions adopted under Chicago's Shared Housing Ordinance change \emph{rental prices} for long-term tenants?
  \item \textbf{Why hard:} treatment is \emph{staggered} -- different tracts adopt in different months -- and effects plausibly vary by cohort and time-since-adoption.
  \item \textbf{Implication for design:} pooled two-way fixed effects (TWFE) implicitly compares newly-treated tracts to \emph{already-treated} ones, generating negative weights and biased ATT under heterogeneity (Goodman-Bacon, 2021).
  \item \textbf{Our spine:} the Callaway--Sant'Anna (CS) group-time ATT estimator with \emph{never-treated} controls; pretrend matched sample; ACS + tract linear trend residualization for robustness.
\end{itemize}
\end{frame}

% --------------------------------------------------------------------
\section{Data and design}

\begin{frame}{Data}
\begin{itemize}
  \item \textbf{Outcome:} ZORI smoothed monthly rent index, ZIP-level, crosswalked to census tract using HUD area shares -- a tract-month panel of rents in dollars.
  \item \textbf{Treatment:} first month a tract has \emph{any} parcel covered by an STR prohibition (Chicago Data Portal: prohibited buildings).
  \item \textbf{Sample:} restricted to tracts whose pre-treatment rent slope can be matched to never-treated comparators (k-NN on slope, $k=3$, $\geq 6$ pre-periods).
  \item \textbf{Result (after Steps 2--5):} matched panel $=$ 373 treated tracts + 114 never-treated; CS keeps 12 prohibition-\emph{cohort months} each with $\geq 5$ tracts (covers 320 of the 373 treated tracts --- very small rollout months are pooled out of aggregated curves).
\end{itemize}
\vfill
{\footnotesize\textit{Estimands operate on the rent index in \$/month; magnitudes are differences in rent levels, not in growth rates.}}
\end{frame}

% --------------------------------------------------------------------
\begin{frame}{Geography of treatment (matched DiD panel)}
\vspace{-0.3em}
\centering
\WPFig{did_story_map.png}{height=0.83\textheight,keepaspectratio}
\end{frame}

% --------------------------------------------------------------------
\section{From raw data to matched sample}

\begin{frame}{Step 1 --- Build one tract--month panel (before matching)}
\vspace{-0.2em}
\textbf{Audience hook: ``What data did we actually start from?''}
\begin{enumerate}\setlength\itemsep{0.45em}
  \item \textbf{Outcome:} smoothed rents (ZORI) measured at ZIP, allocated to census tracts via HUD ZIP--tract weights.
  \item \textbf{Treatment timing:} first calendar month each tract intersects prohibited-building STR registrations (Chicago Data Portal).
  \item \textbf{Panel bookkeeping:} one row per tract per month --- e.g.\ 842 Chicago tracts in the raw file; aligned to 128 months (2015-01 through 2025-08).
\end{enumerate}
\vfill
{\footnotesize Still no causal estimator --- uniform panel structure is prerequisite for every subsequent plot.}
\end{frame}

\begin{frame}{Step 2 --- Why shrink controls? (\emph{slope matching})}
\begin{itemize}\setlength\itemsep{0.4em}
  \item Compared neighbourhoods differ in baseline prices; we nonetheless need untreated tracts whose \emph{pre-prohibition monthly rent trajectory} resembles treated tracts' trajectories as closely as feasible.
  \item \textbf{Procedure:} for each treated tract, retain up to three never-treated tracts with closest pre-period rent slopes ($k{=}3$ NN in slope space); require $\geq 6$ pre months.
  \item \textbf{Treated preservation:} \emph{all} prohibition tracts survive this filter; removals apply only to never-treated tracts that never serve as usable counterfactual slopes.
\end{itemize}
\end{frame}

\begin{frame}{Step 3 --- Tract funnel (counts)}
\vspace{-0.45em}
\centering
\WPFig{data_funnel.png}{height=0.45\textheight,keepaspectratio}\\[-0.1em]
{\scriptsize\setlength{\tabcolsep}{3pt}%
\renewcommand{\arraystretch}{0.92}%
\setlength{\fboxsep}{0pt}%
\resizebox{\linewidth}{!}{\begin{minipage}{\linewidth}%
\WPInput{tables/tab_data_funnel.tex}\end{minipage}}%
}
\vskip0.35em
\footnotesize
\begin{itemize}\setlength\itemsep{0.12em}
  \item Stage B drops many raw never-treated tracts (\emph{never} dropping treated census tracts outright).
  \item Stage C lists CS-visible adoption months --- cohort waves with sufficient mass ($n\geq 5$ tracts).
\end{itemize}
\end{frame}

\begin{frame}{Step 4 --- Descriptive snapshots by prohibition cohort}
\vspace{-0.35em}
\centering
{\scriptsize\setlength{\tabcolsep}{2.5pt}\renewcommand{\arraystretch}{0.95}%
\resizebox{\linewidth}{!}{\begin{minipage}{\linewidth}\WPInput{tables/tab_cohort_descriptive.tex}\end{minipage}}}%
\vskip0.4em
\footnotesize
\begin{itemize}\setlength\itemsep{0.15em}
\item ``Cohort'' $=$ calendar month each tract appears in the prohibition layer for the first time (group-time CS inputs).
\item Early prohibition waves correlate with higher baseline rents on average; income / education columns require a live Census pull (\texttt{CENSUS\_API\_KEY})---rent still sorts cohorts when those fields are blank.
\end{itemize}
\end{frame}

\begin{frame}{Step 5 --- Slopes aligned, levels still tilted (\emph{residualization})}
\begin{columns}[T,totalwidth=\linewidth]
\column{0.53\linewidth}\footnotesize
\begin{itemize}\setlength\itemsep{0.25em}
  \item Matching enforces similarity in \textit{rates of change}; it does \textbf{not} eliminate level wedges (cheap vs.\ expensive neighbourhoods remain cheap vs.\ expensive).
  \item \textbf{Residualized specification:} fit on \emph{pre-treatment months only},
        \[ \widehat Y_{it} \;=\; \hat\alpha + \hat\gamma' X_i^{\text{\tiny ACS}}
              + (\text{tract-specific linear } t) ,\qquad
           \widetilde Y_{it} \;=\; Y_{it} - \widehat Y_{it} + \overline Y. \]
  \item Re-run pooled CS on $\widetilde Y$ (\texttt{drc}) --- magnitudes relocate because subtracting ACS + slopes changes the indexing of dollars.
\end{itemize}
\column{0.46\linewidth}\centering
\scriptsize Pre-window mean rents (months $<$ first prohibition)\\[0.3em]
\renewcommand{\arraystretch}{1.12}
\begin{tabular}{lrrr}
\toprule
\textbf{Group} & Mean & Std & Tracts\\
\midrule
Ever treated & \$1{,}379.67 & 313.86 & 373 \\
Never treated & \$1{,}165.59 & 252.18 & 114 \\
\midrule
\multicolumn{4}{l}{\textit{t}-gap $\approx\$214$,\ $p{<}10^{-11}$ (\texttt{did\_descriptive\_pre\_balance})}\\
\bottomrule
\end{tabular}

\vskip0.55em\raggedright
\textbf{Interpretation for policy slides:} even ``matched'' contrasts mix different rent levels --- hence transparent raw vs.\ residualized CS ladders.
\end{columns}
\end{frame}

% --------------------------------------------------------------------
\begin{frame}{Adoption dynamics and parallel rents}
\begin{columns}[T,onlytextwidth]
\column{0.49\linewidth}
\centering
\WPFig{did_adoption_curve.png}{width=\linewidth}
\column{0.49\linewidth}
\centering
\WPFig{did_parallel_trends.png}{width=\linewidth}
\end{columns}
\vskip0.4em
{\footnotesize Rent paths visually parallel pre-prohibition; the level gap motivates residualization.}
\end{frame}

% --------------------------------------------------------------------
\section{Estimator}

\begin{frame}{Why Callaway--Sant'Anna (and not pooled TWFE)}
\begin{itemize}
  \item For each cohort $g$ and calendar period $t$, the CS estimator forms
        $\widehat{\mathrm{ATT}}(g,t)
         \;=\; \mathbb{E}[Y_t-Y_{g-1}\mid G_g=1]\;-\;\mathbb{E}[Y_t-Y_{g-1}\mid C=1],$
        comparing only to \emph{never-treated} (or not-yet-treated) units.
  \item Aggregation to event-study time $e$: $\widehat{\mathrm{ATT}}(e)=\sum_{g}\omega_{g,e}\,\widehat{\mathrm{ATT}}(g,g+e)$ with positive, transparent weights.
  \item \textbf{What this buys us:} no contamination from already-treated controls, no negative weights, and an interpretable per-cohort dynamic effect.
  \item \textbf{Doubly robust analogue (\texttt{drc}):} residualize $Y$ on pre-period ACS covariates and tract-specific linear trends, then re-run CS. Aligns with Sant'Anna and Zhao (2020).
\end{itemize}
\end{frame}

% --------------------------------------------------------------------
\section{Results}

\begin{frame}{Baseline CS event study (no covariate adjustment)}
\centering
\WPFig{did_callaway_santanna_event_study.png}{width=0.85\linewidth}\\[0.2em]
{\footnotesize Pooled CS ATT estimates on the raw rent index. Pre-period coefficients fluctuate around zero; the post-trajectory drifts mildly negative.}
\end{frame}

% --------------------------------------------------------------------
\begin{frame}{CS event study with ACS + tract-linear-trend controls}
\centering
\WPFig{did_callaway_santanna_event_study_with_controls.png}{width=0.85\linewidth}\\[0.2em]
{\footnotesize Same estimator on \emph{residualized} rent: ACS composition + tract-specific linear trends are netted out before re-running CS. Magnitudes change because the scale changes.}
\end{frame}

% --------------------------------------------------------------------
\begin{frame}{Pooled ATTs: baseline vs.\ residualized CS}
\vfill
\centering
\renewcommand{\arraystretch}{1.2}
\begin{tabular}{lrrr}
\toprule
\textbf{Estimator} & \textbf{ATT (\$/mo.)} & \textbf{SE} & \textbf{95\% CI} \\
\midrule
Callaway--Sant'Anna (matched sample)        & $-2.95$  & 0.63 & $[-4.19,\ -1.71]$ \\
CS w/ ACS + tract linear trends (\texttt{drc}) & $-85.32$ & 0.53 & $[-86.37,\ -84.27]$ \\
\bottomrule
\end{tabular}
\vskip0.7em
\begin{itemize}\footnotesize
  \item Both ATTs are statistically distinguishable from zero; the residualized estimate is interpretable on the residualized scale.
  \item The gap reflects the level wedge documented earlier: residualization sweeps out tract-specific drift that the unconditional CS leaves in.
\end{itemize}
\vfill
\end{frame}

% --------------------------------------------------------------------
\begin{frame}{Cohort heterogeneity (baseline CS)}
\centering
\WPFig{did_cohort_dynamics.png}{height=0.82\textheight,keepaspectratio}
\end{frame}

% --------------------------------------------------------------------
\begin{frame}{Cohort heterogeneity (CS with controls)}
\centering
\WPFig{did_cohort_dynamics_with_controls.png}{height=0.82\textheight,keepaspectratio}
\end{frame}

% --------------------------------------------------------------------
\begin{frame}{Why residualization moves the headline}
\centering
\WPFig{cohort_dynamics_explainer.png}{height=0.78\textheight,keepaspectratio}
\end{frame}

% --------------------------------------------------------------------
\begin{frame}{TWFE vs.\ CS: is the bias guarantee binding here?}
\begin{columns}[T,onlytextwidth]
\column{0.6\linewidth}
\centering
\WPFig{did_twfe_vs_cs_comparison.png}{width=\linewidth}
\column{0.4\linewidth}
\scriptsize
\textbf{Why?}
\begin{itemize}\setlength{\itemsep}{0.25em}
\item TWFE uses \emph{already-treated} units as controls for newly-treated ones (``forbidden comparisons'') $\Rightarrow$ negative weights, biased ATT under heterogeneity.
\item CS only compares to \emph{never-treated} controls -- no negative weights by construction.
\end{itemize}
\textbf{Takeaway}
\begin{itemize}\setlength{\itemsep}{0.25em}
\item Trajectories look parallel; mean TWFE$-$CS gap is \$2.83.
\item TWFE bias is \emph{not} the binding constraint in this Chicago sample, but we still report CS because its guarantee is unconditional.
\end{itemize}
\end{columns}
\end{frame}

% --------------------------------------------------------------------
\section{Robustness: SUTVA / spillovers}

\begin{frame}{Are controls contaminated by neighbour treatment?}
\begin{itemize}
  \item SUTVA failure $\Rightarrow$ never-treated tracts that border prohibition tracts may absorb spillover demand and bias the comparison.
  \item \textbf{Two probes:}
        \begin{enumerate}
            \item \emph{Donut:} restrict the never-treated pool to tracts with \emph{no} ever-treated neighbour, then to those with at least one.
            \item \emph{Dose:} sort never-treated tracts into quartiles of mid-sample share-of-neighbours-treated; re-estimate pooled CS in each.
        \end{enumerate}
\end{itemize}
\end{frame}

% --------------------------------------------------------------------
\begin{frame}{Donut comparison: testing for spillover into controls}
\centering
\WPFig{sutva_donut.png}{height=0.42\textheight,keepaspectratio}\\[0.2em]
\begin{columns}[T,onlytextwidth]
\column{0.48\linewidth}
\scriptsize
\textbf{Idea}
\begin{itemize}\setlength{\itemsep}{0.15em}
\item Neighbours = tracts that share a border (\texttt{geom.touches}, queen contiguity).
\item Split never-treated controls into \emph{isolated} (no treated neighbour) vs \emph{adjacent} ($\geq 1$ treated neighbour).
\item Spillover hypothesis: adjacent controls absorb demand $\Rightarrow$ ATT pulled toward zero.
\end{itemize}
\column{0.48\linewidth}
\scriptsize
\textbf{What we are seeing}
\begin{itemize}\setlength{\itemsep}{0.15em}
\item Full $-\$2.95$; Adjacent $-\$8.18$ ($n{=}100$); Isolated $+\$34.40$ ($n{=}14$).
\item Clean SUTVA failure $\Rightarrow$ isolated should be \emph{more} negative; we see the opposite.
\item 14 isolated tracts are peripheral neighbourhoods $\Rightarrow$ selection, not spillover, drives the gap.
\end{itemize}
\end{columns}
\end{frame}

% --------------------------------------------------------------------
\begin{frame}{Neighbour-treatment dose}
\begin{columns}[T,onlytextwidth]
\column{0.5\linewidth}
\centering
\WPFig{sutva_dose_response.png}{width=\linewidth}
\column{0.5\linewidth}
\scriptsize
\textbf{Construction.} Sort never-treated controls by mid-sample share of treated neighbours; re-run pooled CS in each quartile.
\begin{itemize}\setlength{\itemsep}{0.2em}
\item Q1 (share $0.13$, $n{=}31$): $+\$16.83$.
\item Q2 (share $0.44$, $n{=}32$): $-\$7.71$.
\item Q3 (share $0.66$, $n{=}26$): $-\$15.77$.
\item Q4 (share $0.92$, $n{=}25$): $-\$8.05$.
\end{itemize}
\textbf{Why we are seeing this}
\begin{itemize}\setlength{\itemsep}{0.25em}
\item Q1 mirrors the donut ``isolated'' subset: low-exposure controls sit in spatially distinct neighbourhoods.
\item Q2--Q4 settle into a stable negative band $\Rightarrow$ no monotone dose gradient; pattern reads as neighbourhood selection rather than dose-dependent spillover.
\end{itemize}
\end{columns}
\end{frame}

% --------------------------------------------------------------------
\begin{frame}{Pretrend sensitivity (heuristic)}
\vfill
\centering
{\footnotesize\WPInput{tables/tab_honest_pretrends.tex}}
\vskip0.5em
\begin{itemize}\footnotesize
  \item Reports the maximum pre-period TWFE coefficient and the coarsest discrete second-difference (a smoothness budget) over pre-coefs.
  \item \emph{Spirit only}: a calibrated bound under the Rambachan--Roth (2023) framework would replace these heuristics in a fully ``honest'' sensitivity exercise.
\end{itemize}
\vfill
\end{frame}

% --------------------------------------------------------------------
\section{Takeaways}

\begin{frame}{What we conclude}
\begin{itemize}
  \item \textbf{Tract-level STR prohibitions} are associated with a small \emph{negative} change in tract rents on the unadjusted CS scale and a larger negative effect once tract-specific drift and ACS composition are netted out.
  \item \textbf{Design implication:} reporting both the ``raw'' CS and the residualized CS is honest: residualization changes the meaning of the ATT (it is now on the residual scale), not just its magnitude.
  \item \textbf{Robustness:} cohort heterogeneity is real; spatial probes do not flag obvious SUTVA failure; pretrend coefficients are small in magnitude.
  \item \textbf{Limits:} tract rents are interpolated from ZIP-level ZORI; ``treatment'' is binary even though prohibition density varies; no calibrated R\&R bound is reported here.
\end{itemize}
\end{frame}

% --------------------------------------------------------------------
\begin{frame}{Methods references}
\begin{itemize}\footnotesize
  \item Callaway, B.\ and Sant'Anna, P.\ (2021). \emph{Difference-in-differences with multiple time periods}. Journal of Econometrics 225(2), 200--230.
  \item Goodman-Bacon, A.\ (2021). \emph{Difference-in-differences with variation in treatment timing}. JoE 225(2), 254--277.
  \item Sant'Anna, P.\ and Zhao, J.\ (2020). \emph{Doubly robust difference-in-differences estimators}.
  \item Rambachan, A.\ and Roth, J.\ (2023). \emph{A more credible approach to parallel trends}. RESTUD 90(5), 2551--2598.
  \item Pipeline source: \texttt{src/housing/scripts/did\_pipeline\_callaway\_santanna.py}; documentation in \texttt{docs/CS\_WITH\_CONTROLS.md}.
\end{itemize}
\end{frame}

\end{document}
